\documentclass[11pt,letterpaper, reqno]{article}

\textwidth      =  6in
\textheight     =  8.25in
\oddsidemargin  =  18pt
\evensidemargin =  18pt
\topmargin      =  0.00in

\usepackage[utf8]{inputenc}
\usepackage[spanish]{babel}
\usepackage{amsmath}
\usepackage{amsfonts}
\usepackage{amssymb}
\usepackage{graphicx}
\graphicspath{{img/}}

\title{Análisis de dos Clasificadores}
\author{Abraham Toriz\\ Jesús Mejía\\ Luis Cortés\\ Roberto Saucedo}
\begin{document}
\pagestyle{empty}
\begin{minipage}{.15\textwidth}
	\includegraphics[scale=.22]{uv.eps}
\end{minipage}
\begin{minipage}{.7\textwidth}
\centering
{\huge \sc Universidad Veracruzana}\\
{\LARGE \sc Facultad de Matemáticas}
\end{minipage}
\begin{minipage}{.15\textwidth}
	\begin{flushright}
		\includegraphics[scale=.15]{mateLogo}
	\end{flushright}
\end{minipage}

\vspace{.7cm}
\begin{center}
{\huge \bf
	Análisis de dos Clasificadores
}
\\[1.5cm]
{ \Huge
	T A R E A
}
\\[1.5cm]

{ \large
	Que para pasar la materia
}
\\[.5cm]
{ \LARGE \bf
	Temas Selectos de Matemáticas Aplicadas
}
\\[1.5cm]
{ \Large
I N T E G R A N T E S:
}\\
{ \LARGE \bf
	Abraham Toriz\\
	Jesús Mejía\\
	Luis Cortés\\
	Roberto Saucedo
}\\[1cm]
{ \Large
CATEDRÁTICO:
}\\
{ \LARGE \bf
	Martha Lorena Avendaño
}
\end{center}
\vspace{1.5cm}
\begin{minipage}{.5\textwidth}
\centering
Xalapa, Veracruz
\end{minipage}
\begin{minipage}{.5\textwidth}
\centering
4 de septiembre de 2015
\end{minipage}

%-----------------%
\pagebreak
\pagestyle{plain}

\begin{abstract}
Lorem
\end{abstract}

\section{Regresión Logística}
En este apartado se describirá el modelo de interés...

\section{Los $k$ Vecinos más Cercanos}
En este apartado se describirá el modelo de interés...

\section{Observaciones}

\end{document}